% Section 6: insert after the construction of the two-dimensional limiting scheme.
% This file assumes the paper's existing macros \E, \He, \sgn, \varphi, and \erf.

\section{The Correlation Coefficients}
\label{sec:coefficients}

The previous section defined our two-dimensional limiting scheme, whose two
correlation maps are \(\rho(t)\) and \(t\). We now compute its correlation
function
\[
  H(t)=\sum_{m\geq 1} b_m t^m.
\]
The key observation is that \(g(w,x)=f(w,-x)\). Thus, the Hermite
coefficients of \(f\) determine the entire series \(H\).

We use standard facts about Hermite polynomials. See  \cite{Andrews_Askey_Roy_1999} for a reference.  

Let $\mu$ be the standard Gaussian measure on $\R$, so that
\[
    d\mu(x):=\frac{1}{\sqrt{2\pi}}e^{-x^2/2}\,dx.
\]
Let \(\HermBasis_n\) denote the \(n\)-th orthonormal probabilist's Hermite polynomial. Then \(\{\HermBasis_n\}_{n\ge0}\) is an orthonormal basis for \(L^2(\mu)\).

Let \(W,X\) be independent standard Gaussians. For \(a,b\geq 0\), write
\begin{equation}
  \mathsf A_{a,b}
  =
  \E\!\left[
    f(W,X)\He_a(W)\He_b(X)
  \right].
  \label{eq:mathsfAab}
\end{equation}
Then, in \(L_2\) of the standard Gaussian measure on \(\mathbb{R}^2\),
\[
  f(w,x)
  =
  \sum_{a,b\geq 0}
  \mathsf A_{a,b}\He_a(w)\He_b(x).
\]
Since \(f(-w,-x)=-f(w,x)\), we have \(\mathsf A_{a,b}=0\) whenever
\(a+b\) is even. For a power series \(P\), we write \([t^m]P(t)\) for the
coefficient of \(t^m\).

\begin{theorem}[Coefficient formula]
\label{thm:coefficient-formula}
For every \(t\in(-1,1)\),
\begin{equation}
  H(t)
  =
  \frac{\pi}{2}
  \sum_{a,b\geq 0}
  \mathsf A_{a,b}^2\rho(t)^a(-t)^b.
  \label{eq:H-coefficient-formula}
\end{equation}
Consequently, for every \(m\geq 1\),
\begin{equation}
  b_m
  =
  \frac{\pi}{2}
  \sum_{\substack{a,b\geq 0\\a+b\leq m}}
  (-1)^b\mathsf A_{a,b}^2
  [t^{m-b}]\rho(t)^a.
  \label{eq:bm-coefficient-formula}
\end{equation}
In particular, \(b_m=0\) whenever \(m\) is even.
\end{theorem}

\begin{proof}
Let \((W,X)\) and \((W',X')\) be the Gaussian inputs used to define the
correlation function of the scheme. Thus, \(W\) and \(W'\) have correlation
\(\rho(t)\), while \(X\) and \(X'\) have correlation \(t\); the two
coordinate pairs are independent.

Since \(g(w,x)=f(w,-x)\), the Hermite expansion of \(g\) is
\[
  g(w,x)
  =
  \sum_{a,b\geq 0}
  (-1)^b\mathsf A_{a,b}\He_a(w)\He_b(x).
\]
When the two expansions are substituted into the definition of \(H\),
orthogonality leaves only matching Hermite degrees. The Hermite covariance
identity then contributes \(\rho(t)^a\) in the first coordinate and \(t^b\)
in the second. This gives \eqref{eq:H-coefficient-formula}. The series is
absolutely convergent for \(t\in(-1,1)\), since \(f\in L_2\) and
\(|\rho(t)|\leq |t|<1\).

Extracting the coefficient of \(t^m\) gives
\eqref{eq:bm-coefficient-formula}. Finally, \(\rho\) is odd and
\(\mathsf A_{a,b}\) vanishes unless \(a+b\) is odd. Hence every nonzero
summand in \eqref{eq:H-coefficient-formula} is odd in \(t\), so only odd
coefficients occur.
\end{proof}

\paragraph{Computing the Hermite coefficients.}
The coefficients in \eqref{eq:mathsfAab} reduce to
one-dimensional Gaussian integrals. Let \(Z\) be a standard Gaussian and,
for \(a\geq 0\), set
\begin{equation}
  q_a(u)
  =
  \E\!\left[\sgn(Z+u)\He_a(Z)\right].
  \label{eq:threshold-hermite-coefficients}
\end{equation}
Conditioning on \(X\) gives
\begin{equation}
  \mathsf A_{a,b}
  =
  \E\!\left[
    \He_b(X)q_a\!\left(\vartheta\He_3(X)\right)
  \right].
  \label{eq:f-coefficients-one-dimensional}
\end{equation}
The threshold coefficients have the closed form
\begin{equation}
  q_0(u)
  =
  \erf\!\left(\frac{u}{\sqrt{2}}\right),
  \qquad
  q_a(u)
  =
  \frac{2\varphi(u)}{\sqrt{a}}\He_{a-1}(-u)
  \quad (a\geq 1),
  \label{eq:threshold-hermite-closed-form}
\end{equation}
where \(\varphi(u)=(2\pi)^{-1/2}e^{-u^2/2}\). This standard identity follows
by one integration by parts; we include the proof in
Appendix~\ref{app:coefficient-formulas}.

Theorem~\ref{thm:coefficient-formula} therefore reduces each \(b_m\) to
finite coefficient extraction from powers of \(\rho\) and one-dimensional
Gaussian integrals. This is the formula used to compute the coefficient head
in the next section. A more general formula for schemes built from an
arbitrary finite set of Hermite degrees is recorded in
Appendix~\ref{app:coefficient-formulas}; it is not needed for the main
argument.

